\chapter{Tâche 1 — Framework de test d'intégration et de non-régression}
\label{chap:t1}
% Entreprise : Free Mobile

% =====================================================================
%  Trame conseillée (réutilisée pour chaque tâche) :
%  1. Contexte & dette technique de départ
%  2. Objectif de la mission
%  3. Conception / choix techniques
%  4. Réalisation (extraits de code COMMENTÉS, schémas — PAS de code entier)
%  5. Difficultés rencontrées
%  6. Résultats & bénéfice business (réduction de la dette)
% =====================================================================

\section{Contexte et dette technique initiale}
% TODO : avant ton arrivée, comment la non-régression était-elle gérée ?
% (tests manuels ? absence de filet de sécurité ? régressions récurrentes ?)
% Quantifier si possible le « coût » de cette dette.
\emph{[À rédiger.]}

\section{Objectif de la mission}
% TODO : concevoir un framework de tests d'intégration / non-régression /
% fonctionnels pour le programme initial.
\emph{[À rédiger.]}

\section{Conception et choix techniques}
% TODO : architecture du framework (runner, fixtures, jeux de données,
% rapports), langage/outils retenus et pourquoi, intégration au pipeline.
% Un SCHÉMA d'architecture est le bienvenu :
% \begin{figure}[H]\centering
%   \includegraphics[width=0.8\textwidth]{archi-framework.png}
%   \caption{Architecture du framework de test.}\label{fig:archi-t1}
% \end{figure}
\emph{[À rédiger.]}

\section{Réalisation}
% TODO : illustrer par un EXTRAIT de code commenté (pas le code entier).
% Exemple de gabarit listing :
% \begin{lstlisting}[language=Python,caption={Définition d'un cas de test.},label={lst:t1}]
% def test_exemple():
%     ...
% \end{lstlisting}
\emph{[À rédiger.]}

\section{Difficultés rencontrées}
% TODO : ex. environnements non reproductibles, données de test, flakiness.
\emph{[À rédiger.]}

\section{Résultats et bénéfice business}
% TODO : couverture obtenue, régressions détectées, temps gagné,
% confiance dans les livraisons. Relier au chapitre \ref{chap:veille}.
\emph{[À rédiger.]}
